\documentclass[10pt,a4paper]{amsart}
\usepackage{amsmath,amssymb,amsthm,graphicx}
\usepackage{verbatim}
\newenvironment{reference}{\comment}{\endcomment}
%\newenvironment{reference}{}{}
\newenvironment{slogan}{\comment}{\endcomment}
\newenvironment{history}{\comment}{\endcomment}

% For commutative diagrams we use Xy-pic
\usepackage[all]{xy}

% We use 2cell for 2-commutative diagrams.
\xyoption{2cell}
% Theorem environments.
%
\theoremstyle{plain}
\newtheorem{theorem}[subsection]{Theorem}
\newtheorem{proposition}[subsection]{Proposition}
\newtheorem{lemma}[subsection]{Lemma}

\theoremstyle{definition}
\newtheorem{definition}[subsection]{Definition}
\newtheorem{example}[subsection]{Example}
\newtheorem{exercise}[subsection]{Exercise}
\newtheorem{situation}[subsection]{Situation}

\theoremstyle{remark}
\newtheorem{remark}[subsection]{Remark}
\newtheorem{remarks}[subsection]{Remarks}

\numberwithin{equation}{subsection}

% Macros
%
\def\lim{\mathop{\mathrm{lim}}\nolimits}
\def\colim{\mathop{\mathrm{colim}}\nolimits}
\def\Spec{\mathop{\mathrm{Spec}}}
\def\Hom{\mathop{\mathrm{Hom}}\nolimits}
\def\Ext{\mathop{\mathrm{Ext}}\nolimits}
\def\SheafHom{\mathop{\mathcal{H}\!\mathit{om}}\nolimits}
\def\SheafExt{\mathop{\mathcal{E}\!\mathit{xt}}\nolimits}
\def\Sch{\mathit{Sch}}
\def\Mor{\mathop{\mathrm{Mor}}\nolimits}
\def\Ob{\mathop{\mathrm{Ob}}\nolimits}
\def\Sh{\mathop{\mathit{Sh}}\nolimits}
\def\NL{\mathop{N\!L}\nolimits}
\def\CH{\mathop{\mathrm{CH}}\nolimits}
\def\proetale{{pro\text{-}\acute{e}tale}}
\def\etale{{\acute{e}tale}}
\def\QCoh{\mathit{QCoh}}
\def\Ker{\mathop{\mathrm{Ker}}}
\def\Im{\mathop{\mathrm{Im}}}
\def\Coker{\mathop{\mathrm{Coker}}}
\def\Coim{\mathop{\mathrm{Coim}}}

% Boxtimes
%
\DeclareMathSymbol{\boxtimes}{\mathbin}{AMSa}{"02}

%
% Macros for moduli stacks/spaces
%
\def\QCohstack{\mathcal{QC}\!\mathit{oh}}
\def\Cohstack{\mathcal{C}\!\mathit{oh}}
\def\Spacesstack{\mathcal{S}\!\mathit{paces}}
\def\Quotfunctor{\mathrm{Quot}}
\def\Hilbfunctor{\mathrm{Hilb}}
\def\Curvesstack{\mathcal{C}\!\mathit{urves}}
\def\Polarizedstack{\mathcal{P}\!\mathit{olarized}}
\def\Complexesstack{\mathcal{C}\!\mathit{omplexes}}
% \Pic is the operator that assigns to X its picard group, usage \Pic(X)
% \Picardstack_{X/B} denotes the Picard stack of X over B
% \Picardfunctor_{X/B} denotes the Picard functor of X over B
\def\Pic{\mathop{\mathrm{Pic}}\nolimits}
\def\Picardstack{\mathcal{P}\!\mathit{ic}}
\def\Picardfunctor{\mathrm{Pic}}
\def\Deformationcategory{\mathcal{D}\!\mathit{ef}}
\usepackage{cleveref}
\crefname{originalsection}{original section}{original sections}
\Crefname{originalsection}{Original section}{Original sections}
\makeatletter
\newlabel{example-singleton-sheaf}{{00W3}{source}{}{sourcepointer-1}{}}
\newlabel{example-singleton-sheaf@cref}{{[originalsection][0][]00W3}{[1][0][]source}{}{}{}}
\makeatother
\crefname{originallemma}{original lemma}{original lemmas}
\Crefname{originallemma}{Original lemma}{Original lemmas}
\makeatletter
\newlabel{lemma-special-equivalence}{{03A1}{source}{}{sourcepointer-2}{}}
\newlabel{lemma-special-equivalence@cref}{{[originallemma][0][]03A1}{[1][0][]source}{}{}{}}
\makeatother
\title{Replacing a topos by an equivalent one presented by a good site}
\author{The Stacks Project Authors}
\date{stacks-project:sites.tex:local-snapshot-20261012; GNU FDL 1.2 or later}
\begin{document}
\maketitle
\noindent\textit{Editorial note.} All delivered mathematical text is the author's own, copied from the pinned Stacks Project chapter file sites-cohomology.tex; no statement or proof step was rewritten, reordered or removed.\par\medskip
\noindent\textit{Editorial note.} The proof is the author's complete argument; every step the author writes out is inside the counted span and no closing sentence is a placeholder.\par\medskip
\noindent\textit{Editorial note.} Every labelled reference occurring inside the counted spans is carried in assembly.external\_references as a non-load-bearing navigation pointer (same-chapter ones at their definition site, cross-chapter ones at the invocation inside the counted span), so no load-bearing mathematics is replaced by a stub. No bibliography citation occurs inside the counted spans.\par\medskip
\noindent\textit{Editorial note.} The checked chapter list for same-name collisions is the set of stacks chapters staged in this repository: chow, cohomology, divisors, etale-cohomology, flat, more-algebra, more-morphisms, morphisms, spaces-more-morphisms, varieties together with this card's chapter sites-cohomology.\par\medskip
\noindent\textit{Editorial note.} The preamble is the chapter's own retained minimal preamble (sources/stacks-sites-cohomology/retained-preamble.tex), a verbatim excerpt of the native preamble: the Xy-pic setup and the native theorem declarations with the mathematical macros.\par\medskip
\noindent\textit{Allowed standard prerequisites (editorial).}\begin{itemize}
\item The site attached to a collection of sheaves. Same chapter, Lemma lemma-special-equivalence. Supplies the equivalence of topoi and the properties of the resulting site.
\item The final sheaf. Same chapter, Example example-singleton-sheaf. Supplies the final object included in the first stage of the construction.
\end{itemize}\medskip
\ifdefined\hypertarget\hypertarget{sourcepointer-1}{}\fi
\noindent\textit{Source reference.} Original section 00W3: sources/stacks-sites/pinned-original-source.tex, source lines 7333--7333. Navigation pointer invoked at proof line 7333; the counted span contains the pointer only and the cited result is not reproduced inside it, so the pointer is registered as a non-load-bearing navigation stub in the established form for Stacks material. No bibliography citation is involved.\par\medskip
\ifdefined\hypertarget\hypertarget{sourcepointer-2}{}\fi
\noindent\textit{Source reference.} Original lemma 03A1: sources/stacks-sites/pinned-original-source.tex, source lines 7351--7351. Navigation pointer invoked at proof line 7351; the counted span contains the pointer only and the cited result is not reproduced inside it, so the pointer is registered as a non-load-bearing navigation stub in the established form for Stacks material. No bibliography citation is involved.\par\medskip
\begin{lemma}
\label{lemma-topos-good-site}
Let $\Sh(\mathcal{C})$ be a topos. Let $\{\mathcal{F}_i\}_{i \in I}$
be a set of sheaves on $\mathcal{C}$. There exists an equivalence of topoi
$g : \Sh(\mathcal{C}) \to \Sh(\mathcal{C}')$ induced by a special
cocontinuous functor $u : \mathcal{C} \to \mathcal{C}'$ of sites
such that
\begin{enumerate}
\item $\mathcal{C}'$ has a subcanonical topology,
\item a family $\{V_j \to V\}$ of morphisms of $\mathcal{C}'$
is (combinatorially equivalent to) a covering of $\mathcal{C}'$
if and only if $\coprod h_{V_j} \to h_V$ is surjective,
\item $\mathcal{C}'$ has fibre products and a final object
(i.e., $\mathcal{C}'$ has all finite limits),
\item every subsheaf of a representable sheaf on $\mathcal{C}'$
is representable, and
\item each $g_*\mathcal{F}_i$ is a representable sheaf.
\end{enumerate}
\end{lemma}

\begin{proof}
Consider the full subcategory
$\mathcal{C}_1 \subset \Sh(\mathcal{C})$ consisting of all
$h_U^\#$ for all $U \in \Ob(\mathcal{C})$, the given sheaves
$\mathcal{F}_i$ and the final sheaf $*$ (see
Example \ref{example-singleton-sheaf}). We are going to inductively
define full subcategories
$$
\mathcal{C}_1 \subset \mathcal{C}_2 \subset \mathcal{C}_2 \subset \ldots
\subset \Sh(\mathcal{C})
$$
Namely, given $\mathcal{C}_n$ let $\mathcal{C}_{n + 1}$ be the full
subcategory consisting of all fibre products and subsheaves of objects
of $\mathcal{C}_n$. (Note that $\mathcal{C}_{n + 1}$ has a set
of objects.) Set
$\mathcal{C}' = \bigcup_{n \geq 1} \mathcal{C}_n$.
A covering in $\mathcal{C}'$ is any family
$\{\mathcal{G}_j \to \mathcal{G}\}_{j \in J}$ of morphisms of objects
of $\mathcal{C}'$ such that
$\coprod \mathcal{G}_j \to \mathcal{G}$ is surjective
as a map of sheaves on $\mathcal{C}$.
The functor $v : \mathcal{C} \to \mathcal{C'}$ is given by
$U \mapsto h_U^\#$. Apply
Lemma \ref{lemma-special-equivalence}.
\end{proof}

\end{document}
